\chapter{Build: A Matching Engine and Exchange Simulator}\label{ch:mx:build-a-matching-engine-and-exchange-simulator}

Every tutorial in the books that follow needs a venue: something that accepts orders, matches them, publishes what happened in a feed that can drop
packets, and does exactly the same thing every time it is run with the same inputs. This chapter describes the one this book has used since
chapter~1, \texttt{firm.exchsim}: what an exchange does, in messages; the \omterm{def:mx:the-limit-order-book:engine}{matching engine}; the feed and its recovery; the \omterm{def:mx:build-a-matching-engine-and-exchange-simulator:oe}{order-entry sessions}; and
the clock that makes it deterministic. It then certifies it: the conformance suite, the byte-identity of its three engines, and their speed on this
laptop.

\section{What an exchange does, in messages}

\begin{definition}[Exchange simulator]\label{def:mx:build-a-matching-engine-and-exchange-simulator:sim}
An \emph{exchange simulator}\index{exchange simulator} is a program that behaves toward its participants as a venue does: it accepts and acknowledges
orders through \omterm{def:mx:build-a-matching-engine-and-exchange-simulator:oe}{order-entry sessions}, matches them by the venue's rules, and publishes the resulting \omterm{def:mx:the-limit-order-book:event}{order-book events} in a sequenced market-data feed,
with the latencies, losses and failures of a real connection, and with results that depend only on its inputs.
\end{definition}

\begin{definition}[Order-entry protocol, order-entry session, execution report]\label{def:mx:build-a-matching-engine-and-exchange-simulator:oe}
An \emph{order-entry protocol}\index{order-entry protocol} is the message set a participant uses to send, modify and cancel orders and the venue uses to
answer. An \emph{order-entry session}\index{order-entry session} is one logged-in connection that carries it, with its own sequence of messages, limits and
state. An \emph{execution report}\index{execution report} is the venue's message on that session about one order: accepted, executed (partly or fully),
cancelled, rejected, replaced.
\end{definition}

\texttt{firm.exchsim} speaks four binary protocols, all generated from one table, \texttt{MESSAGES} in the codec module, with a machine-readable copy
in \texttt{schema.json} from which One Quant Book 13 builds its codecs. The \emph{feed} has fourteen message types in an ITCH-like layout (Book~1's add,
execute, cancel, delete and trade messages byte for byte, plus replace, execute-at-price, cross, imbalance, system event, trading action,
stock directory, snapshot header and trailer); order entry has five (new order, replace, cancel, mass cancel, mass quote); \omterm{def:mx:build-a-matching-engine-and-exchange-simulator:oe}{execution reports} six (accepted, executed,
cancelled, rejected, replaced, system); and a control protocol of ten messages drives the venue (system events, logins and drops, phases, bands,
indicatives, crosses, cut-offs, expiries, reference quotes). \texttt{PROTOCOL.md} gives every rule; Nasdaq's TotalView-ITCH~5.0 specification is the model for the
feed (big-endian integers, a stock locate at the same place in every message, nanosecond timestamps since midnight).

A client's day, in process (\texttt{mx\_exchsim.scenario}), shows the reports: a buy of 300 at 100.00 is accepted; a sell of 100 at 100.00 from the
same session trades against it (\omterm{def:mx:build-a-matching-engine-and-exchange-simulator:stp}{self-trade prevention} is off for this session) and both sides get an execution; an immediate-or-cancel sell of 500 at
99.99 executes 200 against the rest and its remaining 300 are cancelled with reason ``I''; a post-only buy of 200 at 100.02 rests; a sell of 100 at
100.01 executes against it at 100.02; a buy at 100.015 is rejected as off the tick grid (``X''); a cancel of the already-filled sell is rejected as too
late (``L''); at the end of the day the post-only order's last 100 shares expire (``E'').

\section{The matching engine}

The engine (\texttt{firm\_exchsim\_engine}) keeps one book per instrument (chapter~1's \texttt{firm.lob}) and applies chapter~2's order types through
\texttt{firm.ordertypes}: limit and \omterm{def:mx:the-limit-order-book:orders}{market orders}, immediate-or-cancel and fill-or-kill, post-only, hidden and \omterm{def:mx:order-types-and-their-uses:hidden}{iceberg orders}, midpoint and primary pegs,
stops, minimum quantities, and \omterm{def:mx:build-a-matching-engine-and-exchange-simulator:stp}{self-trade prevention}.

\begin{definition}[Self-trade prevention, cancel on disconnect]\label{def:mx:build-a-matching-engine-and-exchange-simulator:stp}
\emph{Self-trade prevention}\index{self-trade prevention} stops two orders of the same participant (or group) from trading with each other, by cancelling
the resting order, the incoming one, both, or decrementing them. \emph{Cancel on disconnect}\index{cancel on disconnect} cancels a session's resting
orders when the session is lost.
\end{definition}

Matching gives price priority, then allocates a \omterm{def:mx:the-limit-order-book:tick}{price level} by its rule: first in, first out (price-time priority, Book~1, chapter~19), or pro rata with an
optional top-order share (Book~1's \texttt{firm.match}); it stops at a price band (chapter~25) and at the order's limit.

\omcode{code/firm/exchsim/firm_exchsim_engine.py}{331}{355}{The allocation of one price level in the engine: the eligible resting orders (minimum quantities respected), filled in time order for price-time priority or by Book 1's configurable pro-rata rule, each fill an execution at the resting price.}\label{lst:mx:build-a-matching-engine-and-exchange-simulator:match}

Phases come from a schedule: pre-open and opening call, continuous trading or \omterm{def:mx:auction-mechanics-and-theory:fba}{frequent batch auctions} (chapter~10), closing call and auction, halts and
pauses with reopening calls (chapter~25's \texttt{firm.halts} policies act through the same controls), all uncrossed by chapter~10's auction
(\texttt{firm.auctionsim}). Throttles limit each session's message rate with a token bucket; fees are charged per execution by liquidity flag.

\section{The market-data feed: sequencing, redundancy, recovery}

\begin{definition}[Redundant feed lines, snapshot recovery, retransmission request]\label{def:mx:build-a-matching-engine-and-exchange-simulator:feed}
\emph{Redundant feed lines}\index{redundant feed lines} carry the same sequenced packets twice, on separate paths, so that a receiver can take each packet
from whichever line delivers it first. \emph{Snapshot recovery}\index{snapshot recovery} rebuilds a book from a periodic snapshot of its state, tagged with
the last incremental sequence number it includes, followed by the incremental messages after it. A \emph{retransmission request}\index{retransmission
request} asks the venue's recovery service for messages by sequence number.
\end{definition}

The feed travels in MoldUDP64 packets (a session name, the sequence number of the first message, a message count; a count of zero is a heartbeat) on two
lines, A and B, each with its own impairments drawn from counter-based random streams: independent loss, bursts of loss by a two-state Gilbert--Elliott
channel, duplication, jitter, outages. CME's MDP~3.0 documentation describes the same practice: process both feeds, take each sequence number from
whichever arrives first, and treat a gap on both as loss that needs recovery; its snapshot messages carry the last incremental sequence number processed
so that a receiver can splice the two.

Ten minutes of Book~7's synthetic tape through the venue (\texttt{mx\_exchsim.recovery}), with bursts of loss drawn independently on each line and a
ten-second outage on line B: of 5,793 messages, line A lost 308 and line B 517 (166 of them in the outage); taking each message from either line left 30 missing, and 30 \omterm{def:mx:build-a-matching-engine-and-exchange-simulator:feed}{retransmission requests} filled them,
so the receiver had every message. A snapshot from the middle of the day (the venue took one every 30 seconds) plus the incremental messages after
it rebuilt a book identical, ten levels deep on each side, to the engine's own.

\omcode{code/microstructure/26-build-a-matching-engine-and-exchange-simulator/python/mx_exchsim.py}{118}{123}{Line arbitration's leftovers: the sequence numbers neither line delivered, each asked of the retransmission service, whose MoldUDP64 packet carries the missing message.}\label{lst:mx:build-a-matching-engine-and-exchange-simulator:gaps}

\section{The order-entry protocol and its sessions}

Participants connect through sessions: an agent in process has one per venue with its own latency model (entry, acknowledgement and data delays, jitter,
spikes); a live client connects to the C++20 server on this machine over SoupBinTCP, Nasdaq's point-to-point session protocol over TCP (login with a
requested sequence number, sequenced data, heartbeats), and receives the feed by loopback multicast. Sessions carry the venue's rules: \omterm{def:mx:build-a-matching-engine-and-exchange-simulator:stp}{cancel on
disconnect}, self-trade groups, throttles, and client order identifiers numbered per session (a multi-venue agent must make its own unique across venues,
as chapter~18's router did). \omterm{def:mx:build-a-matching-engine-and-exchange-simulator:oe}{Execution reports} go only to the session that owns the order; the feed goes to everyone, a drop copy to whoever asks.

\section{Clock, latency and determinism}

The simulator is a discrete-event loop on a simulated clock (One Quant Book 7, chapter~17): every message, timer and control is an event at a nanosecond
time, and ties are broken by a fixed order of event classes and insertion. Every random number (latencies, feed impairments, the background's flow) comes
from a counter-based generator (One Quant Book 4, chapter~26) keyed by the seed and by what it is for, so that adding a participant does not change the
random numbers any other one sees. Running the same two-minute day twice gives the same bytes: the SHA-256 of line A's recording and of the engine's
journal agree (\texttt{f7ede7cfa7b3\dots}).

The engine exists three times: the Python reference, a C++20 engine and a Rust engine. A shared fixture (four minutes of tape flow plus a chaos agent
that sends every message type and every control: 17,059 journal records) is replayed by each, and each must reproduce the Python engine's output, feed
messages, reports and final snapshots, byte for byte. Measured on a laptop (Intel Core Ultra 7 155H) under WSL2, with no isolated cores, the Python
engine processes the fixture at about 60,000 records a second; the C++20 test program (which also decodes and checks 4,866 golden messages and compares
every byte) runs through it at more than 500,000 a second. The live C++20 server serves the same engine to clients on localhost.

\section{Tutorial: certification day}\label{tut:mx:build-a-matching-engine-and-exchange-simulator:tut}

\begin{tutorial}
\textbf{Goal.} Trade on the simulator in process and through its server, recover a lossy feed, and certify the three engines.
\textbf{End state:} the numbers of sections 1, 3 and 5.
\begin{tutsteps}
\item \textbf{In process.} \texttt{Simulator(ExchangeConfig)}, an \texttt{Agent} with \texttt{on\_report}; \texttt{mx\_exchsim.scenario()}.
\item \textbf{Feed.} \texttt{FeedConfig(burst\_loss, outages, snapshot\_every\_ns)}; \texttt{res.recorded(line)}, \texttt{res.retransmit(venue, seq, count)},
the snapshots; \texttt{recovery()}.
\item \textbf{Determinism and identity.} \texttt{determinism()}; \texttt{code/firm/exchsim/cpp/exchsim\_engine\_test.cpp} and \texttt{cargo test} in
\texttt{code/firm/exchsim/rust} on the fixture; \texttt{throughput()}.
\item \textbf{Live.} \texttt{code/firm/exchsim/cpp/build.sh}, then \texttt{bin/exchsim\_server --config cfg.json} and a SoupBinTCP client
(\texttt{code/firm/exchsim/tests/test\_serve.py} shows one in Python).
\end{tutsteps}
\textbf{What to change next.} Add implied calendar-spread matching (chapter~21 priced it outside the engine); feed the simulator from chapter~27's agents
as a background; certify a fourth engine against the fixture.
\end{tutorial}

\section{Build: exchange simulator}\label{bld:mx:build-a-matching-engine-and-exchange-simulator:exchsim}

\begin{build}
\bfield{Purpose} The venue under every simulated market of Books 10 to 13: matching, phases and auctions, halts, fees, feeds, sessions.
\bfield{Interface} \texttt{ExchangeConfig(venue, instruments, fees, phases, halts, throttle, engine\_ns, speed\_bump\_ns, batch\_interval\_ns, feed,
faults, reference)}, \texttt{Simulator(venues, sip, seed)} with \texttt{add\_background}, \texttt{add\_agent}, \texttt{add\_events},
\texttt{schedule\_call}, \texttt{run}; \texttt{Result} with \texttt{tape}, \texttt{feed\_messages}, \texttt{recorded}, \texttt{retransmit},
\texttt{snapshot\_bytes}, \texttt{reports}, \texttt{sip}; presets; the C++20 engine, codec and server; the Rust engine and codec (contract:
\texttt{code/firm/INTERFACES.md}, section~7).
\bfield{Rules} Big-endian binary messages generated from one table; prices in 1/10,000 of a currency unit; counter-based randomness keyed by purpose;
no wall clock anywhere in the engine.
\bfield{Acceptance tests} \texttt{code/firm/exchsim/tests/}: thirteen engine-rule scenarios; codec round trips and the schema; Book~1's decoder on the
feed; determinism by hash; the Python engine against the fixture's SHA-256; sequencing, loss and recovery; snapshot rebuild; agents and tape export;
consolidated feed and dark venues; the C++20 and Rust engines byte-identical on the fixture; the Python and C++20 servers round trip.
\bfield{Stretch} Implied spreads; a FIX gateway; a second feed layout in Simple Binary Encoding.
\end{build}

\begin{omsources}
\item Nasdaq, TotalView-ITCH 5.0 specification; MoldUDP64 protocol specification; SoupBinTCP specification.
\item CME Group, MDP~3.0: incremental feed arbitration; market data snapshot full recovery (client systems wiki).
\item FIX Trading Community, Simple Binary Encoding.
\item E.~N.~Gilbert, ``Capacity of a burst-noise channel'', \emph{Bell System Technical Journal} 39(5), 1960.
\end{omsources}

\section{Exercises}

\begin{exercise}[$\star$]\label{exo:mx:build-a-matching-engine-and-exchange-simulator:1}
A MoldUDP64 packet has sequence number 5,001 and a count of 12. What sequence number does the next data packet carry?
\end{exercise}

\begin{exercise}[$\star$]\label{exo:mx:build-a-matching-engine-and-exchange-simulator:2}
From the recovery run, what share of the messages did each line lose, and what share did both lose?
\end{exercise}

\begin{exercise}[$\star$]\label{exo:mx:build-a-matching-engine-and-exchange-simulator:3}
In the scenario, why was the cancel at 7 seconds rejected, and what reason would a cancel of an unknown order get?
\end{exercise}

\begin{exercise}[$\star\star$]\label{exo:mx:build-a-matching-engine-and-exchange-simulator:4}
Why must each random stream be keyed by what it is for, not drawn from one generator in event order?
\end{exercise}

\begin{exercise}[$\star\star$]\label{exo:mx:build-a-matching-engine-and-exchange-simulator:5}
Why does \omterm{def:mx:build-a-matching-engine-and-exchange-simulator:feed}{snapshot recovery} need the snapshot's last sequence number, and what goes wrong without it?
\end{exercise}

\begin{exercise}[$\star\star$]\label{exo:mx:build-a-matching-engine-and-exchange-simulator:6}
What does byte-identity of three engines on a fixture prove, and what does it not?
\end{exercise}

\begin{exercise}[$\star\star\star$]\label{exo:mx:build-a-matching-engine-and-exchange-simulator:7}
\emph{Coding.} Run the recovery with line B's outage but no burst loss on either line. How many messages does each line lose, and how many
\omterm{def:mx:build-a-matching-engine-and-exchange-simulator:feed}{retransmission requests} are needed?
\end{exercise}

\begin{exercise}[$\star\star\star$]\label{exo:mx:build-a-matching-engine-and-exchange-simulator:8}
\emph{Find the flaw.} ``The C++ engine is eight times faster than the Python one, so the simulator's sessions will run eight times faster in C++.''
\end{exercise}

\section{Problem: Certification Day}

\begin{problem}[{Weekend problem --- certification day}]\label{pb:mx:build-a-matching-engine-and-exchange-simulator:1}
Before other teams rely on the simulator, it must be certified: rules, recovery, determinism, identity across implementations and speed.

\medskip\noindent\textbf{Part I --- Messages.}
\begin{enumerate}
\item Define an \omterm{def:mx:build-a-matching-engine-and-exchange-simulator:sim}{exchange simulator}, an \omterm{def:mx:build-a-matching-engine-and-exchange-simulator:oe}{order-entry protocol}, a session and an \omterm{def:mx:build-a-matching-engine-and-exchange-simulator:oe}{execution report}.
\item List the four protocols and their message counts.
\item Walk through the scenario's reports.
\item Define \omterm{def:mx:build-a-matching-engine-and-exchange-simulator:stp}{self-trade prevention} and \omterm{def:mx:build-a-matching-engine-and-exchange-simulator:stp}{cancel on disconnect}.
\end{enumerate}

\medskip\noindent\textbf{Part II --- The engine.}
\begin{enumerate}[resume]
\item How does the engine allocate a \omterm{def:mx:the-limit-order-book:tick}{price level} under each rule?
\item What stops matching?
\item Which phases can the venue be in, and what uncrosses a call?
\item How do \texttt{firm.halts} policies act on the venue?
\end{enumerate}

\medskip\noindent\textbf{Part III --- The feed.}
\begin{enumerate}[resume]
\item Describe a MoldUDP64 packet and the two lines.
\item How are losses and outages drawn?
\item Give the recovery run's losses, gaps after arbitration and retransmissions.
\item How was the book rebuilt from a snapshot, and how was it checked?
\item What does CME's MDP~3.0 documentation recommend for its A and B feeds?
\end{enumerate}

\medskip\noindent\textbf{Part IV --- Certification.}
\begin{enumerate}[resume]
\item What makes the simulator deterministic?
\item Describe the fixture.
\item \emph{State the named result}: the conformance suite passed by each engine, the byte-identity of the three engines' feeds on the fixture, and the
engines' messages per second on this laptop.
\item Why is the C++ number a lower bound on the engine's speed?
\item Which parts of the simulator are not certified by the fixture?
\item What would you add before a live trading team used it?
\item In one sentence: why does a simulator need to be deterministic?
\end{enumerate}
\end{problem}

\section{Interview questions}

\begin{interviewq}[$\star$ \iqroles{developer}]\label{iq:mx:build-a-matching-engine-and-exchange-simulator:1}
How does a feed handler recover from a gap in a UDP multicast feed?
\end{interviewq}

\begin{interviewq}[$\star\star$ \iqroles{developer}]\label{iq:mx:build-a-matching-engine-and-exchange-simulator:2}
Design a \omterm{def:mx:the-limit-order-book:engine}{matching engine}'s data structures for price-time priority with fast cancels.
\end{interviewq}

\begin{interviewq}[$\star\star$ \iqroles{developer}]\label{iq:mx:build-a-matching-engine-and-exchange-simulator:3}
How would you make a market simulator reproducible to the byte, including its randomness?
\end{interviewq}

\begin{interviewq}[$\star\star$ \iqroles{trader}]\label{iq:mx:build-a-matching-engine-and-exchange-simulator:4}
Your IOC order was partly filled and the rest cancelled. What reports do you receive, in what order?
\end{interviewq}

\begin{interviewq}[$\star\star$ \iqroles{risk}]\label{iq:mx:build-a-matching-engine-and-exchange-simulator:5}
Why do venues offer \omterm{def:mx:build-a-matching-engine-and-exchange-simulator:stp}{cancel on disconnect} and \omterm{def:mx:build-a-matching-engine-and-exchange-simulator:stp}{self-trade prevention}, and what risks does each control?
\end{interviewq}

\begin{interviewq}[$\star\star\star$ \iqroles{developer}]\label{iq:mx:build-a-matching-engine-and-exchange-simulator:6}
Two implementations of the same engine disagree on one message out of a million. How do you find out which is wrong?
\end{interviewq}
